\clearpage
\heading{What the plainness check establishes}\label{sec:15evidence}
\smallref{Pass 15 verification. \pub{NS1/Accel.lean}{27}{Definition of plainness}, \pub{Paper/Postpone.lean}{30}{K.4} and \pub{Paper/NoAcc.lean}{216}{conditional substitution}.}

\subhead{A check in the actual machine}
Consider the program that quotes the empty tree and then uses the result as the name of a scalar. Quoting costs one step. The name instruction costs one more and appends the empty word $\epsilon$ to its bill. On the empty input its sole accepting record has cost two, ordinary weight one and bill $[\epsilon]$. The list $[\epsilon]$ has length one. Thus the program halts, meets a two-step accepting budget, and fails plainness at every time allowance. This example lives in the publisher's language for every choice of the base weights $\theta$.

\subhead{An independent composition check}
Let $p$ halt on $z$, let $T<2^J$ bound its accepting costs, and let the name lookup coincide with the tally trie's lookup on every request the program actually makes. Assume that K.4's displayed upper bounds for the size and time of each wrapper fit the accelerator's input and plainness budgets. Write $\kappa$ for the original name lookup and $K.\mathrm{lk}$ for the lookup used to value the accelerator. Then
\[
 \operatorname{val}_{\kappa}(p,z)
  =\sum_e\lambda^e\,\operatorname{val}_{K.\mathrm{lk}}(D,z_e).
\]
To prove it, expand the left side by K.4. For the current summand take the size inequality and plainness derivation supplied by that same application of K.4. Compose size inequalities, enlarge the plain time allowance, and use the value clause at that input. Summing over $e$ proves the formula. This isolates the substitution in K.6. Establishing its hypotheses for the self-program is the work on the preceding pages.

The two lemmas above and the three tally-endpoint lemmas on p.~\pageref{sec:15endpoint} were checked with Lean 4.33.1 after rebuilding the 16 publisher modules imported by \texttt{Paper.NoAcc}, counting that target. The targeted run read back the types and axioms of eight publisher declarations and all five review lemmas. Its external library objects were prebuilt at the release revisions. The full project and the independent second kernel have their earlier execution records on p.~\pageref{sec:14formal}. The finite Python controls exercise the boundary examples and selected resource schedules.

\subhead{An optional route through the implementation}
K.5 consumes one cutting program with bounds on all short descriptions. Trace J.3's encoded outcome through J.4's iterated list and recover, for the same cut run, its cost, list length and piece sizes. Then recover piece plainness and the weighted sum of products under the parent's plainness assumption. This would check the producer of the precise contract used on p.~\pageref{sec:15k5}.

\smallref{Repository files include \texttt{checks/pass15\_plain.lean}, \texttt{checks/pass15\_lean.json}, and the three per-track controls. \pub{Paper/Cuts.lean}{460}{J.4 program} · \pub{NS1/Scale.lean}{73}{K.5 consumer}.}
